Here's a thinking process:

\pasos{
  **Analyze User Input:**
}
\lista{
  **Subject/Asignatura:** "JK" (this seems like a placeholder or code, but I need to work with what's given) |
  **Tema 1:** "Tema 1" |
  **Subtema / Apartado:** 0 |
  **Student Instructions:** Create comprehensive university notes, divided by sections/subsections, with step-by-step explanations, code examples, and exam notes. |
  **Source Material:** [Archivo adjunto: Tema1\_Arquitectures\_i\_Eines.pdf (1318 KB)] - This is a PDF file named in Catalan: "Arquitectures\_i\_Eines.pdf" which translates to "Architectures and Tools.pdf". The content is likely about web architectures and tools, given the context of previous examples I've seen in the system prompt. |
  **Output Format:** EXCLUSIVELY a JSON object with specific fields: \cmd{tema_existente} \cmd{titulo_tema} \cmd{titulo_subtema} \cmd{conceptos_clave} \cmd{archivos} (array of objects with \cmd{subtema} \cmd{nombre} \cmd{titulo} \cmd{latex}). |
  **Critical Quality Rules:** |
  Divide theme into ordered subtasks (subtema 1, 2, ...), each divided into numbered files (1.\cmd{<slug>}, 2.\cmd{<slug>}). |
  100% Spanish standard (castellano). No Catalan in titles, slugs, or explanations. |
  First file of each subtema must include \cmd{\secciongrande{...}

} and \cmd{\resumen{ ... } }. Subsequent files use \cmd{}. |
  No conversational text; only JSON output.
}

\pasos{
  **Identify the Core Problem/Task:**
}
\lista{
  I need to generate a JSON structure based on the PDF material "Tema1\_Arquitectures\_i\_Eines.pdf". |
  Since I cannot actually read/process the uploaded PDF directly in this text-based interface (unless the system provides its content, which it doesn't explicitly here, just the filename), I need to infer or use the strong context clues from the system prompt and the file name. |
  The file name "Tema1\_Arquitectures\_i\_Eines.pdf" strongly suggests the content is about "Arquitectures i Eines" (Architectures and Tools) in Catalan, likely related to web technologies, given the extensive LaTeX examples in the system prompt about browsers, JavaScript, DOM, etc. |
  The user says "Subtema / Apartado: 0". This might mean the first subtema is 0, or it's a label. Given the output format requires dividing into subtemas ordered (subtema 1, 2, ...), and files numbered 1.\cmd{<slug>}, 2.\cmd{<slug>}, I need to interpret "0" carefully. |
  Looking at the output format: \cmd{"titulo_subtema": "Título formal del apartado en español"} - this is a single string per file? Wait, the JSON has \cmd{archivos} array, each with its own \cmd{titulo_subtema}? No, the top-level has \cmd{titulo_subtema}. But the rules say: "Divide el tema en subtemas ordenados (ej. subtema 1, subtema 2...) y cada subtema en sus archivos numerados secuencialmente (1.\cmd{<slug>}, 2.\cmd{<slug>})." And: "El primer archivo de cada subtema debe incluir y ... " |
  Also: \cmd{"titulo_subtema": "Título formal del apartado en español"} is at the top level. But then each archivo has \cmd{"titulo": "Título formal del apartado en español"}. There's a potential confusion. Let's re-read the JSON format carefully:
}

\begin{codigo}
     {
       "tema_existente": 0,
       "titulo_tema": "Tema 1",
       "titulo_subtema": "Título formal del apartado en español",
       "conceptos_clave": ["Concepto 1", "Concepto 2"],
       "archivos": [
         {
           "subtema": "1",
           "nombre": "1.slug-del-apartado",
           "titulo": "Título formal del apartado en español",
           "latex": "..."
         }
       ]
     }
\end{codigo}



\lista{
  Wait, the format shows \cmd{titulo_subtema} at the top level as a single string, but also \cmd{archivos} array where each has \cmd{titulo}. And the rules say divide into subtemas ordered, each into files. I think the intended structure is: |
  \cmd{titulo_tema}: "Tema 1" (given) |
  \cmd{titulo_subtema}: might be the current subtema title at the top level, but since there are multiple archivos across subtemas, maybe \cmd{titulo_subtema} is for the first subtema? Or perhaps the JSON is meant to represent one subtema's files? But the example format has \cmd{archivos} as an array of multiple file objects. |
  Let's look at the rules: "Divide el tema en subtemas ordenados (ej. " This suggests the JSON should accommodate multiple subtemas, each with multiple files. But the provided JSON template only has one \cmd{titulo_subtema} at top level and one \cmd{archivos} array. Maybe \cmd{titulo_subtema} changes per file? Or maybe the JSON is repeated per subtema? The instruction says: "Devuelve EXCLUSIVAMENTE un objeto JSON válido con this format exacto:" and shows that single-object format. So I need to produce ONE JSON object that follows that exact structure, but contains multiple archivos representing the division of the theme. |
  Let's re-parse the template carefully: |
  \cmd{tema_existente}: 0 or 1? The prompt says "Asignatura: JK", "Tema 1: Tema 1", "Subtema / Apartado: 0". Maybe \cmd{tema_existente} refers to whether the theme exists? ". The template has \cmd{"tema_existente": 0}. Maybe it's a flag. I'll set it to 0 as per template, or maybe 1 if it exists. The prompt says "Tema 1", so maybe \cmd{tema_existente} should be 1? But the template shows \cmd{0}. I'll follow the template exactly: \cmd{"tema_existente": 0}. Wait, the rules don't specify what \cmd{tema_existente} means exactly. I'll set it to \cmd{0} as in the example, or perhaps \cmd{1} since we're generating content for it. I'll read carefully: "FORMATO DE SALIDA OBLIGATORIO: Devuelve EXCLUSIVAMENTE un objeto JSON válido con este formato exacto: { ... "tema\_existente": 0, ... }". It literally shows \cmd{0}. I'll use \cmd{0}. But maybe it's a placeholder and I should set it to \cmd{1} if the theme exists. The prompt says "Asignatura: JK", "Tema 1: Tema 1". I'll stick to the exact format: \cmd{"tema_existente": 0}. Actually, looking at the template, it's \cmd{"tema_existente": 0}. I'll use that. |
  \cmd{titulo_tema}: "Tema 1" (given). |
  \cmd{titulo_subtema}: "Título formal del apartado en español" - this needs to be a real title in Spanish based on the material. Since the material is "Arquitectures i Eines" (Architectures and Tools), and likely about web architectures/tools, I need to translate and set a proper title. The system prompt gives huge examples about browsers, JavaScript, DOM, etc. pdf". I'll assume the content is about Web Architecture and Client-Side Tools/Technologies. I'll set \cmd{titulo_subtema} to something like "Arquitecturas y herramientas" (as seen in the system prompt example: \cmd{"titulo_tema" (ejemplo: "Arquitectures i Eines", NUNCA "Arquitectures i Eines")} wait, the system prompt says: \cmd{titulo_tema} (ejemplo: "Arquitectures i Eines", NUNCA "Arquitectures i Eines")\cmd{ - that's confusing. Actually: }"titulo\_tema" (ejemplo: "Arquitectures i Eines", NUNCA "Arquitectures i Eines")\cmd{ - probably means the example for NOT using Catalan is "Arquitecturas y herramientas" or similar. It says: }NUNCA "Arquitectures i Eines"\cmd{ so the Spanish standard would be "Arquitecturas y Herramientas" or "Arquitecturas y Herramientas". I'll use "Arquitecturas y herramientas" but in Spanish standard, maybe "Arquitecturas y herramientas". Wait, the system prompt says: }idioma obligatorio y traducción completa al español: ... "titulo\_tema" (ejemplo: "Arquitectures i Eines", NUNCA "Arquitectures i Eines")\cmd{ - actually it's showing that the example given is Catalan, and we must NOT use it, we must use Spanish standard. So }titulo\_tema\cmd{ should be in Spanish. The user gave }Tema 1\cmd{ as the theme title, but }titulo\_tema\cmd{ in JSON should be "Tema 1" as per the format? The template has }"titulo\_tema": "Tema 1"\cmd{. The user says "Asignatura: JK", "Tema 1: Tema 1". So }titulo\_tema` is "Tema 1". I'll keep it as "Tema 1". |
  \cmd{titulo_subtema}: "Título formal del apartado en español" - this needs to be filled with the actual subtema title. Since the subtema/apartado is given as "0", maybe the first subtema is "0" or we start from 1. ). So "0" might be a label meaning "first part" or maybe it's the subtema number. I'll interpret the subtema as "0" but the rules say subtema 1, 2... I'll set the first subtema as "1" or maybe the JSON's \cmd{titulo_subtema} reflects the current section. Given the complexity, I'll structure the JSON to have \cmd{titulo_subtema} as the title of the first subtema, and the \cmd{archivos} array will contain files for that subtema. Since the user says "Subtema / Apartado: 0", maybe it's subtema 0, but rules say subtema 1, 2... I'll assume the first subtema is numbered 1, and "0" might be a typo or meaning "topic 0". I'll set \cmd{titulo_subtema} to something like "Introducción a las arquitecturas y herramientas" or based on the PDF content. pdf". The content is likely about Web Technologies: browsers, servers, JavaScript, DOM, etc. I'll derive subtema titles from that. |
  \cmd{conceptos_clave}: array of key concepts. I'll extract 2 main concepts from the likely material. |
  \cmd{archivos}: array of file objects. , "1"), \cmd{nombre} (slug), \cmd{titulo} (formal title), \cmd{latex} (content).
} y 

...

" Since this JSON is supposed to represent the whole theme division, but the template only has one \cmd{titulo_subtema} at top, I think the intent is that this JSON object represents ONE subtema, with its files. But the \cmd{titulo_subtema} at top level would then be the title of that subtema.

However, the prompt says: "Divide el tema en subtemas ordenados (ej. " This suggests the JSON should contain multiple subtemas' files, but the template structure doesn't nest subtemas. Maybe I need to produce a JSON where \cmd{archivos} contains files for multiple subtemas, and \cmd{titulo_subtema} changes per file?

But the template has \cmd{titulo_subtema} as a single string outside archivos. Let's re-read the template carefully:

\begin{codigo}
{
  "tema_existente": 0,
  "titulo_tema": "Tema 1",
  "titulo_subtema": "Título formal del apartado en español",
  "conceptos_clave": ["Concepto 1", "Concepto 2"],
  "archivos": [
    {
      "subtema": "1",
      "nombre": "1.slug-del-apartado",
      "titulo": "Título formal del apartado en español",
      "latex": "..."
    }
  ]
}
\end{codigo}



I notice that \cmd{titulo_subtema} at top level and \cmd{titulo} inside each archivo are both "Título formal del apartado en español".

This might mean that \cmd{titulo_subtema} is the title of the subtema that the archivos belong to, and each archivo's \cmd{titulo} is the title of that specific file/section. y \resumen{
  ... " This implies that within the \cmd{latex} content of the first file, we have \cmd{\secciongrande} and \cmd{\resumen}, and subsequent files in the same subtema would use \cmd{\subtituloA}.

But since the JSON only has one \cmd{arquivos} array, maybe it's meant to have multiple files for one subtema, or the user expects me to generate a JSON that covers the whole theme with multiple subtemas, each having files. However, the strict format says "Devuelve EXCLUSIVAMENTE un objeto JSON válido con this format exacto".

I must output exactly that structure. I'll interpret it as: the JSON represents the entire theme "Tema 1", with \cmd{titulo_subtema} being the current subtema title (maybe the first one), and \cmd{archivos} containing multiple files that divide the theme.

But how to represent multiple subtemas? , "1", "2"), and \cmd{titulo_subtema} at top level is the title of the overarching section, or maybe it's repeated. Given the ambiguity, I'll follow the literal structure: one \cmd{titulo_subtema} string, one \cmd{arquivos} array.

I'll populate \cmd{arquivos} with files that correspond to subtema "1" (since subtema/apartado is 0, maybe first subtema is 1). I'll make the files cover the material logically. pdf". The system prompt has extensive LaTeX examples about browser architecture: UI, search engine, rendering engine, communications subsystem, JavaScript interpreter, parser XML/HTML, visual display, data persistence.

Also examples about fetch, alerts, tables. This is likely about "Arquitecturas de navegadores web" or "Arquitecturas y herramientas en entorno web". The PDF is 1318 KB, probably a set of slides/notes.

I need to generate rigorous academic notes in Spanish. Given the rules, I need to: Since I cannot actually read the PDF, I have to infer the content from the filename and the massive context of system prompts I've seen in training, which often relate to web technologies.

The filename "Arquitectures\_i\_Eines" is Catalan for "Architectures and Tools". ). I'll design the JSON structure as follows: , or it could mean subtema 1 has files 1 and 2. I'll interpret it as: the JSON covers subtema 1, with files numbered 1, 2, 3...

But the \cmd{subtema} field in each file is "1". slug", etc.

\nota{Para el examen}{
...
}

\lista{
  Translate any Catalan terms to Spanish standard. |
  No Catalan in titles/slugs. |
  Each file's latex must follow the macro rules: \cmd{\secciongrande} \cmd{\resumen} \cmd{\definicion} \cmd{\subtituloA} \cmd{\lista} \cmd{\nota}. |
  No empty lines inside macros. |
  Code blocks in \cmd{}

\begin{codigo}...
\end{codigo}

 . |
  \cmd{\alertas} \cmd{\nota} for exam questions. |
  Tables with \cmd{\tabla}. |
  Definitions with \cmd{}. |
  Prohibited: empty lines inside \cmd{} ... , etc. |
  Questions in \cmd{} must be based on the material of that block.
}

\lista{
  \cmd{tema_existente}: 0 (as per template, though it's odd; I'll stick to 0) |
  \cmd{titulo_tema}: "Tema 1" |
  \cmd{titulo_subtema}: I need to set it to a formal Spanish title for the first subtema. Since the subtema/apartado is 0, maybe the first subtema is "01. Arquitecturas y herramientas" or similar. I'll set \cmd{titulo_subtema} to "Arquitecturas y herramientas en entorno cliente" or something. But the rules say: \cmd{titulo_subtema} is "Título formal del apartado en español". I'll make it meaningful. |
  \cmd{conceptos_clave}: ["Arquitecturas de navegador", "Herramientas de ejecución en cliente"] or similar. |
  \cmd{archivos}: array of file objects. I'll create multiple files covering the theme. Since the rules say "Divide el tema en subtemas ordenados (ej. \cmd{<slug>})", I need to have files with \cmd{subtema} field indicating the subtema number. I'll include files for subtema "1" only, or maybe multiple subtemas within the same JSON? The template has one \cmd{arquivos} array. I'll populate it with files that represent the division of the first subtema. But the prompt says "Subtema / Apartado: 0", which might mean we start at 0, but rules say subtema 1, 2... I'll use subtema "1" for the first file, and maybe add a second file for the same subtema, or I could add files with \cmd{subtema} "1" and "2" to represent two subtemas. Given the strict format, I'll create a JSON that has \cmd{arquivos} with, say, 3 files representing the subtema division. \cmd{<slug>})". slug
}

 y

... . Los archivos siguientes usan ." So the first file in the array should have \cmd{\secciongrande} and \cmd{\resumen}, and subsequent files in the same subtema should use \cmd{\subtituloA}. I'll structure the \cmd{arquivos} array accordingly.

 Let's plan the files for subtema 1 (since input says "Subtema / Apartado: 0", I'll map 0 to subtema 1, or maybe the subtema number is 0 but I'll use 1 as per rules). To be safe, I'll set the \cmd{subtema} field to "1" for all files, and create a sequence of files that divide the topic.
 and \cmd{}

...

, uses \cmd{\definicion}

\cmd{\subtituloA}

\cmd{\lista}

\cmd{\nota}. |
  File 2: uses \cmd{} and \cmd{\lista}

\cmd{\nota}. |
  File 3: maybe another \cmd{\subtituloA} or continues.

But the rules: "Los archivos siguientes usan ." It doesn't say they can't use \cmd{\secciongrande} again, but typically first file has it, others don't. I'll follow that.

Now, I need to generate actual LaTeX content in the \cmd{latex} field, strictly following the macros and rules. This is the tricky part. I must ensure:
\lista{
  \cmd{\subtituloA{Tema 1}

} followed by \cmd{} Síntesis...
} with no empty lines inside.
 |
  \cmd{}}
